\par\bigskip\noindent\textbf{K101 supporting proof: The directed finite chart and period assembly}\par\nobreak
Working ledger owned by blind\_control, saved before checkpoint24.
The full finite quotient resolution is owned by root. Fourfold owns
the source/Miller/point-line atlas EI.1-\/-12. The detailed root and
annular maps are EIR.1-\/-11; the independently checked local toric
maps are SIT.1-\/-5. No numerical return is assigned here.

\par

\medskip\noindent\textbf{ASI.1. The actual domains and overlap arrows}

\par\nobreak

Use the original source parameter and selected branch
\[
 r=e^2(3e^2-2)/(1-3e^2)^2,\qquad
 t_{c,o}=e^4(1-e^2)/(1-3e^2)^3 .
\]
The original finite cap is kept on its actual fixed infinity disc
\(|q_{\infty,o}|<\delta\), with the complete K96 \(A,B\), character
and reference cochain. The ordinary source/Cartier family outside
the collision extends on compact subsets of the original minimal
source: its selected height is a nonzero analytic function there.
Its entire value stays in the lift \(B_{\rm sel}\).

The finite-native source and the complementary ordinary chart have
the exact overlap
\[
 (x_e,y_e,t_e)=(\chi^{-2}\bar X,\chi^{-3}\bar Y,\chi^{-4}),
 \quad e=\chi\sigma,\quad q_e=\chi^{-4}H.
\]
The point frames agree, and their actual line chains satisfy
\[
 X_j=\chi^j A_j,\quad Y_j=\chi^{-j-1}B_j,\qquad
 \lambda e^4\widehat B_e
       =\lambda\sigma^4H\,\mathcal B_{a,b}.
\]
These equalities are EI.74-\/-76 and EIR.54-\/-56, with all original
Miller factors and the genuine point Hom frame included. They join
the finite-native \(I_2\) source resolution to the complementary
ordinary chart. No division section is used at \(\sigma=0\).

On the selected \(H\)-disc, the six-factor chart has its actual
negative \(T_0\) boundary. Its overlap with the genuine ordinary
point-line chain is
\[
 k=4j+r-1,\qquad
 X_k=\psi^{-1}q^j s_r,\quad Y_k=\psi q^{-j}t_r,
 \qquad \psi=f_{P^a}/r_a .
\]
The full \(\psi\), the integer \(q^j\), the source-chain overlap,
the physical \(B\) reset and the remaining fourth-root action
are retained in SIT.11-\/-17. The source section graph is carried
by the nine-factor \((\eta^4,v^4,q)\) atlas with its exact
\((\eta,v)\mapsto(\eta c_b,v/c_b)\) source reset.

The complementary chart reaches the original finite cap only on
its large-\(\sigma\) overlap:
\[
 q_{\infty,o}=d^3/\sigma^4,\qquad
 |\sigma|>|d|^{3/4}\delta^{-1/4}.
\]
There the map is the complete original ordinary-to-cap map, pulled
through the displayed source and point-line graph. Its sector
square is EIR.58-\/-60:
\[
 \Phi_\beta=A^kB^{j_{\alpha\beta}}\Phi_\alpha,\qquad A^m=B^n.
\]
Thus the graph generator \(\sigma\mapsto-i\sigma\), taken along
the stated negative quarter turn, acts on root representatives
by \(AB^{j_\gamma}\). Every finite character, period column and
reference suspension integer stays in this map. The fourth
iterate is \(A^4B^{\sum j_\gamma}\); it is the identity deck
word on the completed cap quotient. This is the exact finite
groupoid descent on the overlap.

\par

\medskip\noindent\textbf{ASI.2. The punctured family identification}

\par\nobreak

For \(e\ne0\), the base graph is an isomorphism: its new charts
have the written inverse source rescalings, and its exceptional
base locus is in \(e=0\). The O/T blowup centres have
\(\sigma=0\), hence lie in \(e=0\). In the selected factor
atlas, the four \(\sigma\) factors are units when \(e\ne0\);
deleting them gives exactly the two original \(T_0\) ruling
charts, with the full unit maps of SIT.11. In the nine-factor
section atlas, \(\eta v=\sigma\ne0\), so its eight unit factors
delete by the same literal chart maps. Thus these modifications
retain the original punctured family.

The actual family map on the ordinary charts is the original
source rescaling, carrying its canonical point section. On the
original finite cap it is the identity in the retained original
cap coordinates. The overlap equality is the preceding
\(\Phi_\beta\) square, including its \(B\)-powers. On the selected
Tate chart the invariant coordinates \(Q,Z,V_{\rm pt}\) agree;
the tensor coordinates carry
\[
 V_{{\rm tensor},o}=(e/d)^{1-a}V_{{\rm tensor},e}.
\]
Every Cox character consequently carries the exact factor
\[
 e^{4h+(1-a)j}d^{-3h+(a-1)j}
\]
from EIR.7. This factor is an invertible analytic function on
the punctured parameter domain, with its written integer
winding retained. No contraction of that winding is used.

In a continued native root coordinate the same map has
\(J_{u+1}=D_o^{-4}J_u\). The complete cap representative is
the compensating \(A^{-4}D^4\), and in the reduced row the
retained \(B^{-n}\) reindex and the full period column of
EIR.28 are included. Alternatively, the preceding construction
uses the fixed original cap coordinates throughout. The
EIR.58-\/-60 square proves that these two descriptions give the
same quotient map, with the actual inverse and all overlap
cochains.

Consequently the punctured atlas carries the original oriented
family and its positive \(e\)-loop through explicit source,
line, cap and selected-ruling maps. This identifies its mapping
torus with the original marked-loop mapping torus. It does not
turn the unresolved central compact assembly into an evaluated
characteristic number.

\par

\medskip\noindent\textbf{ASI.3. Remaining finite normal and compact assembly data}

\par\nobreak

The ordinary O/T section corners retain the two-cone
\(\mu_4\) actions of SIT.25. Root is resolving these in their
actual quotient lattice. The entire chain admits the common
physical-lattice dictionary
\[
 V_{\rm phys}=(\eta v)^jX_j,\qquad
 (x_\chi,x_\eta,x_v,x_X)\longmapsto
 (x_\chi,x_\eta,x_v,x_X+j(x_\eta+x_v)).
\]
It sends
\[
 \mathbb Z^4+\mathbb Z(1,-1,0,j)/4
 \quad\hbox{onto}\quad
 \mathbb Z^4+\mathbb Z(1,-1,0,0)/4
\]
by the displayed integral coordinate change. The physical
circle is the last primitive coordinate. In this common
lattice the two cones at index \(j\) are
\[
 (e_\chi,R^v_j,R^\eta_{j+1},R^\eta_j),\qquad
 (e_\chi,R^\eta_{j+1},R^v_{j+1},R^v_j),
\quad
 R^\eta_j=(0,1,0,j),\ R^v_j=(0,0,1,j).
\]
The actual \(B\) reset is \(j\mapsto j+4\), or
\(x_V\mapsto x_V+4(x_\eta+x_v)\), together with its full
source and Hom units. This gives the required period arrows
for compatible finite-quotient resolutions.

Still to be assembled are their compact fundamental domains,
the complete boundary list, and the normal and physical
circle-line characteristic comparison after all those pieces
are glued. The local source graph contains its actual A1
normalization and the contracted intermediate curve; neither
is removed from that remaining boundary calculation.

\par

\medskip\noindent\textbf{ASI.4. Completed local interfaces after checkpoint24}

\par\nobreak

The factor-order comparison is now the actual common
modification FO.1--6 in
\path{notes/part3_e6_selected_infinity_factor_overlay_working.md.}
It retains the alternating ordinary order \((\eta,v)^4\)
and the selected grouped order \((\eta^4,v^4,q)\).
Their exact overlap has 38 regular three-cones.
The ordinary quotient extension has 238 regular
four-cones in the original invariant lattice; the
full selected extension has 39 regular cones through
\(q=0\). Both parent morphisms, all period faces
and every unit exponent are saved.

In particular, the common dense frame square is
\[
 V_\theta=\psi V,\qquad
 \beta=\mathcal B\,\frac{\psi(T_bp)}{\psi(p)},\qquad
 q=\kappa H .
 \tag{ASI.1}
\]
The ordinary period is
\(V\mapsto\lambda H\mathcal B\,\sigma^4V\).
The selected period is
\(V_\theta\mapsto(\lambda\beta/\kappa)\sigma^4qV_\theta\).
Their equality follows by substitution in ASI.1,
with all arguments retained. Its period-index arrow
is \(q^j\psi^{-1}\) on the \(q\)-unit overlap.
Through \(q=0\), that same \(q^j\) is part of the
literal selected monomial period
\((\eta,v,q,V)\mapsto(\eta,v,q,V+4\eta+4v+q)\).

The 238-cone common quotient has the actual
\(B\)-matrix
\[
 \begin{pmatrix}1&0&0&0\\0&1&0&0\\0&0&1&0\\-1&4&4&1\end{pmatrix}
\]
in \((\chi^4,\chi\eta,v,V)\)-valuation coordinates.
It replaces the ordinary-only 52-cone template
when the common selected interface is required.
The earlier template, all its data and its audit
remain preserved. There are two analytic source
copies, O and T; their period faces glue crosswise
under the full \(B\) map, with \(c_b\) and
\(u_O/u_T\). The cone number is a local chart count.

The native finite I2 graph now includes its entire
\(t=0\) fibre through IA.3--13. Its source ideal
is exactly \((b_0,\epsilon)^2\), and its O
ideal is \((\eta_e,\epsilon)\). On \(\chi\ne0\),
\[
 (b_0)_e=\chi\bar Y,\qquad
 \eta_e=\chi\bar\eta,\qquad\epsilon=\chi\sigma.
\]
The two graph ideals differ from the complementary
ones by the same source unit \(\chi\).
Their projective coordinates, actual point
Cartier ratios and original-source resolutions
therefore agree. The line-chain arrow is the
complete \(X_k=\chi^kA_k,\ 
Y_k=\chi^{-k-1}B_k\).

\par

\medskip\noindent\textbf{ASI.5. The remaining finite source copies}

\par\nobreak

The other two central two-torsion points are
\((\bar X,\bar Y)=(i,0),(-i,0)\).
The fourth-root generator exchanges their
neighbourhoods; its square fixes each.
Their four free coordinates are
\((\chi,\zeta=\bar Y,X_j,Y_j)\), with full
fourth-root weights \((1,-1,j,-j-1)\).
For even \(j\) the actual fixed ideal is
\((\chi,\zeta,Y_j)\); for odd \(j\) it is
\((\chi,\zeta,X_j)\).

The complete source map is
\[
 \bar X'=\frac d{\bar X}
           +\frac{2\sigma(\zeta+\sigma)}{\bar X^2},\quad
 \zeta'=-\frac{\bar X'}{\bar X}\zeta
            +\left(1-\frac{\bar X'}{\bar X}\right)\sigma .
 \tag{ASI.2}
\]
Since \(\sigma=X_jY_j\), this preserves the
actual fixed ideal by the invertible triangular
matrix with its full source-mixing term.
Its inverse preserves the same ideal.
Blowing up this ideal and then taking its
involution quotient gives the three smooth
charts
\((\rho=t^2,z_b,z_c,u)\), where \(u\) is
the neutral coordinate and \(w_i=tz_i\)
are the odd coordinates. Every coarse
quadratic has the literal map
\(w_iw_k=\rho z_i z_k\).
For a full ideal matrix \(M\), put
\(L=Mz\); its target-pivot chart is
\[
 \rho'=\rho L_\beta^2,\qquad
 z_\gamma'=L_\gamma/L_\beta ,
 \tag{ASI.3}
\]
with the entire neutral-line multiplier.
Thus the nonmonomial source mixing extends
through the actual quotient resolution.

The full \(B\) exchanges the two source
points and shifts \(j\) by four. When one
point is chosen to represent the exchanged
pair, its period closing is
\(\gamma^{-1}B\), including the exact
fourth-root phases and physical normal
weights. The adjacent chain-index
identifications retain the nonzero
coordinate squared in the odd-line
ideal. These maps are the inputs to
the global fixed-face ledger.

The marked original volume uses
\[
 -8(X_jY_j)^3d^{-2}u\,
        d\chi\wedge d\zeta\wedge dX_j\wedge dY_j,
 \qquad u=(3\bar X^2+2a_2\bar X+d)^{-1}.
 \tag{ASI.4}
\]
Its central value is \(u=-1/2\) at both
points. The quotient chart\textquotesingle s oriented
full unit is \(-4d^{-2}u\), and its
logarithmic character is \((1,1,4,4)\)
in the four cover coordinates.
The exact source-differential identity
transports this full unit, not just
its central value.

The independent \ensuremath{\mu}4 and \ensuremath{\mu}2 audits are
\path{records/blind_audit_e6_selected_infinity_mu4_quotient.json}
and \path{records/blind_audit_e6_selected_infinity_mu2_pair.json.}
The common factor fan supplies an additional actual
modification with both parent arrows. These local
calculations do not assign the compact normal
double, determinant-line extension or final return;
those require the complete attached fixed-face
and cap ledger.

